\documentclass[10pt,letterpaper]{article}

\usepackage[utf8]{inputenc}
\usepackage[T1]{fontenc}
\usepackage[spanish,es-nodecimaldot]{babel}
\usepackage[landscape,margin=0.25in]{geometry}
\usepackage{amsmath,amssymb}
\usepackage{xcolor}
\usepackage{multicol}
\usepackage{enumitem}
\usepackage{array,tabularx}

\definecolor{azul}{HTML}{173B63}
\definecolor{azulclaro}{HTML}{E8F0F7}
\definecolor{rojo}{HTML}{9A2C2C}
\definecolor{gris}{HTML}{454B52}

\pagestyle{empty}
\setlength{\parindent}{0pt}
\setlength{\parskip}{0.65pt}
\setlength{\columnsep}{0.20in}
\setlength{\columnseprule}{0.25pt}
\setlength{\abovedisplayskip}{1.4pt}
\setlength{\belowdisplayskip}{1.4pt}
\setlength{\abovedisplayshortskip}{0.8pt}
\setlength{\belowdisplayshortskip}{0.8pt}
\setlist[itemize]{leftmargin=9pt,itemsep=0.3pt,topsep=0.7pt,parsep=0pt}
\setlist[enumerate]{leftmargin=12pt,itemsep=0.3pt,topsep=0.7pt,parsep=0pt}

\newcommand{\Htriple}[3]{\{#1\}\,#2\,\{#3\}}
\newcommand{\topic}[1]{%
  \vspace{1.0pt}\colorbox{azulclaro}{%
    \parbox{\dimexpr\linewidth-2\fboxsep\relax}{\color{azul}\bfseries #1}}\vspace{0.5pt}}
\newcommand{\alerta}[1]{\textcolor{rojo}{\bfseries #1}}
\newcommand{\prof}{\textsuperscript{\scriptsize Prof.}}

\begin{document}
\fontsize{6.85}{7.55}\selectfont

\begin{center}
  {\fontsize{13.2}{14}\selectfont\bfseries\color{azul} FORMULARIO DE LÓGICA DE HOARE --- Q1}\\[-1pt]
  {\small Programación Avanzada \(\cdot\) Prof. Gustavo Márquez Flores \(\cdot\) Notas 4--6}
\end{center}
\vspace{-5pt}

\begin{multicols*}{3}

\topic{1. Lectura y notación del profesor}

\[
\boxed{\Htriple{P}{C}{Q}}
\]
Si el estado inicial cumple la \textbf{precondición} $P$ y el código $C$
\emph{termina}, el estado final cumple la \textbf{poscondición} $Q$.

\begin{tabularx}{\linewidth}{@{}>{\bfseries}l X@{}}
$\{P\},\{Q\}$ & conjuntos de relaciones entre variables y constantes (aserciones).\\
$\{\ \}$ & aserción vacía: no restringe nada; equivale a $\top$.\\
$\bot$ & contradicción; ningún estado la satisface.\\
$x_0$ & valor inicial/histórico; no cambia aunque cambie $x$.\\
$:=$ & asignación; $=$ es igualdad lógica.\\
$s\models P$ & el estado $s$ satisface $P$.\\
\end{tabularx}

\textbf{Corrección parcial:} si termina, entrega $Q$.\quad
\textbf{Corrección total:} parcial $+$ prueba de terminación.

\textbf{Validez vs. demostración:}
$\models\Htriple{P}{C}{Q}$ significa ``es verdadera''; 
$\vdash\Htriple{P}{C}{Q}$ significa ``se deriva con las reglas''.

\topic{2. Fuerza lógica y consecuencia}

$R$ es \textbf{más fuerte} que $S$ si $R\Rightarrow S$: admite menos estados.
$\bot$ es la más fuerte; $\top$ la más débil.

\[
\boxed{
\frac{P'\Rightarrow P\qquad \Htriple{P}{C}{Q}\qquad Q\Rightarrow Q'}
     {\Htriple{P'}{C}{Q'}}}
\]
Se puede \textbf{fortalecer la precondición} y \textbf{debilitar la
poscondición}. Mnemotecnia: las flechas apuntan hacia la terna nueva,
$P'\Rightarrow P$ y $Q\Rightarrow Q'$.

\topic{3. Axioma de asignación: regla clave}

Para $V:=E$, partir de la meta $Q$ y sustituir en ella el valor que tendrá $V$:
\[
\boxed{\Htriple{Q_E^V}{V:=E}{Q}},\qquad
Q_E^V\equiv Q[E/V].
\]
Lectura: ``$Q$ con $V$ sustituida por $E$''. En las diapositivas,
$\{V=E,Q\}$ \textbf{indica reemplazo}; \alerta{no} significa $V=E\land Q$.
La condición calculada es la más débil que garantiza $Q$.

\textbf{Ejemplos relámpago}
\[
\begin{aligned}
\Htriple{?}{k:=4a}{k=12}:&\quad 4a=12\iff\boxed{a=3},\\
\Htriple{?}{i:=2i}{i<6}:&\quad 2i<6\iff\boxed{i<3},\\
\Htriple{?}{x:=1/x}{x\ge0}:&\quad 1/x\ge0,\ x\ne0\iff\boxed{x>0}.
\end{aligned}
\]
\alerta{Dominio:} añadir $x\ne0$, radicando $\ge0$, índice válido, etc.

\topic{4. Composición y cálculo hacia atrás}

\[
\boxed{
\frac{\Htriple{P}{C_1}{R}\qquad\Htriple{R}{C_2}{Q}}
     {\Htriple{P}{C_1;C_2}{Q}}}
\]
$R$ es la aserción intermedia. El código se ejecuta hacia delante, pero se
\textbf{verifica de derecha a izquierda}:
\[
Q\xleftarrow{\ C_n\ }R_{n-1}\xleftarrow{\ C_{n-1}\ }\cdots
\xleftarrow{\ C_1\ }P_{\min}.
\]
Al final comprobar $P_{\rm dada}\Rightarrow P_{\min}$.

\textbf{Ejemplo:}
\[
\Htriple{?}{a:=a+2;\ b:=2a}{b=14}
\]
$b=14\Leftarrow a=7\Leftarrow a+2=7$; luego $\boxed{a=5}$.

\topic{5. Reglas que combinan ternas}

Para el \textbf{mismo código} $C$:
\[
\frac{\Htriple{P_1}{C}{Q_1}\quad\Htriple{P_2}{C}{Q_2}}
     {\Htriple{P_1\land P_2}{C}{Q_1\land Q_2}}
\quad
\frac{\Htriple{P_1}{C}{Q_1}\quad\Htriple{P_2}{C}{Q_2}}
     {\Htriple{P_1\lor P_2}{C}{Q_1\lor Q_2}}.
\]
Conjunción acumula garantías; disyunción reúne casos alternativos.
$\Htriple{P}{\mathbf{skip}}{P}$ deja el estado intacto.

\topic{6. Condicionales: demostrar todos los caminos}

\textbf{Con \texttt{else}:}
\[
\boxed{
\frac{\Htriple{P\land B}{C_1}{Q}\qquad
      \Htriple{P\land\neg B}{C_2}{Q}}
     {\Htriple{P}{\mathbf{if}\ B\ \mathbf{then}\ C_1\ \mathbf{else}\ C_2}{Q}}}
\]
Ambas ramas deben garantizar la \textbf{misma} $Q$ (o algo que la implique).

\textbf{Sin \texttt{else}:}
\[
\frac{\Htriple{P\land B}{C_1}{Q}\qquad(P\land\neg B)\Rightarrow Q}
     {\Htriple{P}{\mathbf{if}\ B\ \mathbf{then}\ C_1}{Q}}.
\]
La rama falsa ejecuta $\mathbf{skip}$; no se puede ignorar.

\textbf{Ejemplo del profesor: máximo}
\[
\Htriple{\top}{\mathbf{if}\ a>b\ \mathbf{then}\ m:=a\ \mathbf{else}\ m:=b}
{m\ge a\land m\ge b}.
\]
Rama $B$: sustituir $m\mapsto a$ y usar $a>b\Rightarrow a\ge b$.
Rama $\neg B$: sustituir $m\mapsto b$ y usar $\neg(a>b)\iff b\ge a$.
Si $P\land B=\bot$, esa rama es inalcanzable y su terna vale vacíamente.

\topic{7. Invariantes de bloque y de ciclo}

Una relación $I$ es preservada por el bloque $C$ si
\[
\boxed{\Htriple{I}{C}{I}}.
\]
Las variables pueden cambiar; lo que permanece es la relación.

Para $\mathbf{while}\ B\ \mathbf{do}\ C$:
\[
\boxed{
\frac{\Htriple{I\land B}{C}{I}}
     {\Htriple{I}{\mathbf{while}\ B\ \mathbf{do}\ C}{I\land\neg B}}}
\]
Para demostrar $\Htriple{P}{\mathbf{while}\ B\ \mathbf{do}\ C}{Q}$ hay tres
\textbf{condiciones de verificación}:
\[
\boxed{P\Rightarrow I}\quad
\boxed{\Htriple{I\land B}{C}{I}}\quad
\boxed{I\land\neg B\Rightarrow Q}.
\]
Son, respectivamente: \textbf{inicialización, mantenimiento y salida útil}.

\textbf{Conexión con inducción:}\par
\begin{tabularx}{\linewidth}{@{}l X@{}}
Caso base & $I$ vale antes de la iteración $0$.\\
Hipótesis & $I$ vale después de $n$ iteraciones.\\
Paso & una ejecución de $C$ da $I$ para $n+1$.\\
Salida & usar además $\neg B$ para obtener $Q$.\\
\end{tabularx}

\topic{8. Cómo inventar y fortalecer $I$}

\begin{enumerate}
  \item Nombrar valores iniciales que cambian: $M_0,x_0,\ldots$
  \item Trazar $0,1,2,3$ iteraciones y escribir variables en función de $n$.
  \item Eliminar $n$ y buscar: \emph{acumulado $\circ$ pendiente $=$ total}.
  \item Añadir cotas que permitan concluir $Q$ al salir.
  \item Probar; una tabla sugiere $I$, pero \alerta{no lo demuestra}.
\end{enumerate}

\textbf{Patrones del profesor}\par
\begin{tabularx}{\linewidth}{@{}l X@{}}
Cuadrado por sumas & $I:C=DA$, $V=A-D$.\\
Potencia & $I:RN^M=N^{M_0}$, $V=M$.\\
Suma $1\ldots n$ & $I:s=\frac{i(i+1)}2\land0\le i\le n$, $V=n-i$.\\
\end{tabularx}
Ejemplo de fortalecimiento: $Y=X!$ puede preservarse, pero añadir $X\le N$
permite deducir $X=N$ cuando la salida solo da $X\ge N$.

\topic{9. Terminación: función variante}

El invariante prueba corrección parcial. Para corrección total, encontrar una
variante entera $V$ tal que, mientras $I\land B$:
\[
\boxed{V\ge0}\qquad\text{y}\qquad\boxed{V_{\rm después}<V_{\rm antes}}.
\]
Forma formal útil, guardando el valor anterior $v$:
\[
\Htriple{I\land B\land V=v}{C}{I\land0\le V<v}.
\]
Un entero no negativo no puede disminuir estrictamente para siempre.
\textbf{Cuidado:} preservar $I$ no prueba que el ciclo termine.

\topic{10. Precondición más débil y transformadores}

Para corrección parcial:
\[
\Htriple{P}{C}{Q}\quad\Longleftrightarrow\quad P\Rightarrow wlp(C,Q).
\]
Para código determinista sin ciclos, $wp$ y $wlp$ se calculan igual:
\[
\begin{aligned}
wlp(V:=E,Q)&=Q[E/V],\\
wlp(C_1;C_2,Q)&=wlp(C_1,wlp(C_2,Q)),\\
wlp(\mathbf{if},Q)&=(B\Rightarrow wlp(C_1,Q))\\
&\quad\land(\neg B\Rightarrow wlp(C_2,Q)).
\end{aligned}
\]
$wlp$: precondición liberal más débil (parcial). $wp$: precondición más débil
(total). $sp(C,P)$ razona hacia delante y produce la poscondición más fuerte.

\topic{11. Receta de examen}

\begin{enumerate}
  \item Escribir claramente $P$, código $C$, meta $Q$ y dominios.
  \item Marcar variables modificadas; separar valor actual de inicial ($x_0$).
  \item Sin ciclo: empezar en $Q$ y sustituir \textbf{hacia atrás}.
  \item En secuencias: insertar cada aserción intermedia $R$.
  \item En un \texttt{if}: abrir casos $P\land B$ y $P\land\neg B$.
  \item En un \texttt{while}: proponer $I$ y probar entrada, preservación y salida.
  \item Si piden corrección total: añadir variante y sus dos obligaciones.
  \item Simplificar las fórmulas; usar consecuencia para cerrar las implicaciones.
  \item Si una flecha parece falsa, buscar un contraejemplo concreto.
\end{enumerate}

\topic{12. Errores que cuestan puntos}

\begin{itemize}
  \item Sustituir en $P$ hacia delante en vez de sustituir $V\mapsto E$ en $Q$.
  \item Tratar $\{V=E,Q\}$ como conjunción o confundir $:=$ con $=$.
  \item Omitir el dominio de una operación parcial.
  \item No verificar $P_{\rm dada}\Rightarrow P_{\min}$.
  \item Probar una sola rama del \texttt{if}; olvidar el $\mathbf{skip}$ sin \texttt{else}.
  \item Llamar invariante a algo que solo coincide en ejemplos.
  \item Usar $I$ sin $\neg B$ para concluir la poscondición.
  \item Confundir corrección parcial con total; olvidar la variante.
\end{itemize}

\topic{13. Pruebas de cordura de 20 segundos}

\begin{itemize}
  \item \textbf{Asignación:} toma un valor que cumpla la precondición calculada,
        ejecuta una vez y revisa $Q$.
  \item \textbf{Consecuencia:} una precondición más fuerte admite \emph{menos}
        entradas; una poscondición más débil exige \emph{menos} a la salida.
  \item \textbf{Ciclo:} revisa estado inicial, una vuelta simbólica y estado de salida.
  \item \textbf{Variante:} calcula $V$ antes/después; debe bajar y no cruzar sin límite.
\end{itemize}

\vfill
{\scriptsize\color{gris}
Fuente del curso: \emph{Programación Avanzada, Notas 4--6}, Gustavo Márquez
Flores. Convención usada: $Q_E^V\equiv Q[E/V]$ y $\{\ \}=\top$.
}

\end{multicols*}
\end{document}
